\appendix
\chapter{Documentația API REST Completă}
\label{cap:api_documentation}

Acest capitol prezintă specificația completă a API-ului REST pentru backend-ul StyleGenAI. Documentația include sintaxa requestului, response schemas, coduri de eroare și exemple practice.

\section{Convenții și Standarde}

\subsection{Protocol și Format}

\begin{itemize}
    \item \textbf{Protocol}: HTTP/1.1 cu suport HTTPS (recomanda în production)
    \item \textbf{Content-Type}: \texttt{application/json} (implicit) sau \texttt{multipart/form-data} (pentru upload fișiere)
    \item \textbf{Base URL}: \texttt{http://localhost:5000} (dev) sau \texttt{https://api.stylegenai.com} (production)
    \item \textbf{Encoding}: UTF-8
\end{itemize}

\subsection{Autentificare}

Fiecare request la endpoint-uri protejate trebuie să includă header:

\begin{lstlisting}[language=HTTP]
Authorization: Bearer <access_token>
\end{lstlisting}

unde \texttt{<access\_token>} este JWT token primit din login. Token-ul expiră după 24 ore.

\subsection{Format Răspuns de Eroare}

Erorile sunt returnate în format JSON consistent:

\begin{lstlisting}[language=JSON]
{
  "success": false,
  "error": "Descriere error",
  "error_code": "ERROR_CODE_CONSTANT",
  "status": 400,
  "timestamp": "2026-05-21T10:30:00Z"
}
\end{lstlisting}

---

\section{Autentificare}

\subsection{POST /auth/register}

\textbf{Descriere}: Înregistrare utilizator nou și emitere JWT token.

\textbf{Request}:

\begin{lstlisting}[language=HTTP]
POST /auth/register HTTP/1.1
Content-Type: application/json
\end{lstlisting}

\begin{lstlisting}[language=JSON]
{
  "email": "user@example.com",
  "password": "SecurePassword123!",
  "username": "john_doe"
}
\end{lstlisting}

\textbf{Validații}:
\begin{itemize}
    \item Email: format RFC 5322, unique în DB
    \item Password: minimum 8 caractere, cel puțin 1 uppercase, 1 lowercase, 1 digit, 1 special char
    \item Username: 3-30 caractere, alphanumeric + underscore
\end{itemize}

\textbf{Response Succes (201)}:

\begin{lstlisting}[language=JSON]
{
  "success": true,
  "user_id": 42,
  "email": "user@example.com",
  "username": "john_doe",
  "access_token": "eyJhbGciOiJIUzI1NiIsInR5cCI6IkpXVCJ9...",
  "expires_in": 86400,
  "message": "Utilizator înregistrat cu succes"
}
\end{lstlisting}

\textbf{Response Erori}:

\begin{lstlisting}[language=JSON]
// 400 - Email deja înregistrat
{
  "success": false,
  "error": "Email deja înregistrat",
  "error_code": "EMAIL_EXISTS",
  "status": 400
}
\end{lstlisting}

\begin{lstlisting}[language=JSON]
// 400 - Password slab
{
  "success": false,
  "error": "Password nu întrunește criterii de securitate",
  "error_code": "WEAK_PASSWORD",
  "status": 400,
  "details": "Necesită: 8+ chars, uppercase, lowercase, digit, special char"
}
\end{lstlisting}

---

\subsection{POST /auth/login}

\textbf{Descriere}: Autentificare și emitere token JWT.

\textbf{Request}:

\begin{lstlisting}[language=HTTP]
POST /auth/login HTTP/1.1
Content-Type: application/json
\end{lstlisting}

\begin{lstlisting}[language=JSON]
{
  "email": "user@example.com",
  "password": "SecurePassword123!"
}
\end{lstlisting}

\textbf{Response Succes (200)}:

\begin{lstlisting}[language=JSON]
{
  "success": true,
  "user_id": 42,
  "email": "user@example.com",
  "username": "john_doe",
  "access_token": "eyJhbGciOiJIUzI1NiIsInR5cCI6IkpXVCJ9...",
  "expires_in": 86400
}
\end{lstlisting}

\textbf{Response Erori}:

\begin{lstlisting}[language=JSON]
// 401 - Invalid credentials
{
  "success": false,
  "error": "Email sau parolă incorectă",
  "error_code": "INVALID_CREDENTIALS",
  "status": 401
}
\end{lstlisting}

---

\subsection{POST /auth/refresh}

\textbf{Descriere}: Reînnoire token JWT expirat.

\textbf{Request}:

\begin{lstlisting}[language=HTTP]
POST /auth/refresh HTTP/1.1
Authorization: Bearer <expired_token>
\end{lstlisting}

\textbf{Response Succes (200)}:

\begin{lstlisting}[language=JSON]
{
  "success": true,
  "access_token": "eyJhbGciOiJIUzI1NiIsInR5cCI6IkpXVCJ9...",
  "expires_in": 86400
}
\end{lstlisting}

---

\section{Gestionare Haine (Clothes Management)}

\subsection{POST /clothes/upload}

\textbf{Descriere}: Încărcare articol vestimentar cu imagine.

\textbf{Request}:

\begin{lstlisting}[language=HTTP]
POST /clothes/upload HTTP/1.1
Authorization: Bearer <token>
Content-Type: multipart/form-data
\end{lstlisting}

FormData:
\begin{itemize}
    \item \texttt{image} (File, required): JPG/PNG/WebP, max 10MB
    \item \texttt{category} (String, required): "Top" | "Bottom" | "Shoe"
    \item \texttt{style} (String, optional): "Casual" | "Business" | "Sport" | "Formal"
    \item \texttt{color\_hint} (String, optional): "Red" | "Blue" | etc. (pentru validare)
\end{itemize}

\textbf{Response Succes (201)}:

\begin{lstlisting}[language=JSON]
{
  "success": true,
  "cloth_id": 1001,
  "user_id": 42,
  "category": "Top",
  "style": "Casual",
  "dominant_color": "#FF5733",
  "dominant_color_name": "Red",
  "image_path": "/uploads/user_42/cloth_1001.jpg",
  "image_thumbnail": "/uploads/user_42/cloth_1001_thumb.jpg",
  "created_at": "2026-05-21T10:30:00Z"
}
\end{lstlisting}

\textbf{Response Erori}:

\begin{lstlisting}[language=JSON]
// 413 - File prea mare
{
  "success": false,
  "error": "Fișier prea mare (max 10MB)",
  "error_code": "FILE_TOO_LARGE",
  "status": 413
}
\end{lstlisting}

\begin{lstlisting}[language=JSON]
// 400 - Tip fișier invalid
{
  "success": false,
  "error": "Tip fișier nepermis. Acceptate: JPG, PNG, WebP",
  "error_code": "INVALID_FILE_TYPE",
  "status": 400
}
\end{lstlisting}

---

\subsection{GET /clothes}

\textbf{Descriere}: Listare haine utilizatorului curent.

\textbf{Request}:

\begin{lstlisting}[language=HTTP]
GET /clothes?category=Top&style=Casual&limit=50&offset=0 HTTP/1.1
Authorization: Bearer <token>
\end{lstlisting}

Query Parameters (optional):
\begin{itemize}
    \item \texttt{category}: "Top" | "Bottom" | "Shoe"
    \item \texttt{style}: "Casual" | "Business" | "Sport" | "Formal"
    \item \texttt{limit}: default 50
    \item \texttt{offset}: default 0
    \item \texttt{sort}: "created_at\_desc" (default) | "created_at_asc" | "category"
\end{itemize}

\textbf{Response Succes (200)}:

\begin{lstlisting}[language=JSON]
{
  "success": true,
  "clothes": [
    {
      "cloth_id": 1001,
      "category": "Top",
      "style": "Casual",
      "dominant_color": "#FF5733",
      "image_path": "/uploads/user_42/cloth_1001.jpg",
      "created_at": "2026-05-21T10:30:00Z"
    },
    {
      "cloth_id": 1002,
      "category": "Bottom",
      "style": "Casual",
      "dominant_color": "#2E4053",
      "image_path": "/uploads/user_42/cloth_1002.jpg",
      "created_at": "2026-05-21T10:25:00Z"
    }
  ],
  "total_count": 42,
  "page": 1,
  "per_page": 50
}
\end{lstlisting}

---

\subsection{DELETE /clothes/{cloth\_id}}

\textbf{Descriere}: Ștergere articol vestimentar.

\textbf{Request}:

\begin{lstlisting}[language=HTTP]
DELETE /clothes/1001 HTTP/1.1
Authorization: Bearer <token>
\end{lstlisting}

\textbf{Response Succes (200)}:

\begin{lstlisting}[language=JSON]
{
  "success": true,
  "message": "Articol șters cu succes",
  "cloth_id": 1001
}
\end{lstlisting}

\textbf{Response Erori}:

\begin{lstlisting}[language=JSON]
// 404 - Articol nu găsit
{
  "success": false,
  "error": "Articol nu găsit",
  "error_code": "CLOTH_NOT_FOUND",
  "status": 404
}
\end{lstlisting}

\begin{lstlisting}[language=JSON]
// 403 - Nu este owner
{
  "success": false,
  "error": "Nu ai permisiune să ștergi acest articol",
  "error_code": "FORBIDDEN",
  "status": 403
}
\end{lstlisting}

---

\section{Generare Recomandări}

\subsection{POST /get\_suggestion}

\textbf{Descriere}: Generare outfit recommendation pe baza preferințelor și stării garderobei utilizatorului.

\textbf{Request}:

\begin{lstlisting}[language=HTTP]
POST /get_suggestion HTTP/1.1
Authorization: Bearer <token>
Content-Type: application/json
\end{lstlisting}

\begin{lstlisting}[language=JSON]
{
  "season": "spring",
  "occasion": "casual",
  "filters": [
    "exclude_formal",
    "prefer_bold_colors"
  ],
  "limit_combinations": 1,
  "preferences": {
    "color_harmony": "complementary",
    "coverage": "full"
  }
}
\end{lstlisting}

\textbf{Request Field Descriptions}:

\begin{itemize}
    \item \textbf{season} (String, required): "spring" | "summer" | "autumn" | "winter" - affects trend boost
    \item \textbf{occasion} (String, optional): "casual" | "business" | "sport" | "formal" - filters articole
    \item \textbf{filters} (Array, optional): list de filtre apply-ate
    \item \textbf{limit\_combinations} (Int, optional, default 1): câte ținute to return (1-10)
    \item \textbf{preferences} (Object, optional): personalization hints
\end{itemize}

\textbf{Response Succes (200)}:

\begin{lstlisting}[language=JSON]
{
  "success": true,
  "outfit": {
    "outfit_id": 5001,
    "user_id": 42,
    "components": {
      "top": {
        "cloth_id": 1001,
        "category": "Top",
        "dominant_color": "#FFFFFF",
        "dominant_color_name": "White",
        "image_path": "/uploads/user_42/cloth_1001.jpg",
        "style": "Casual"
      },
      "bottom": {
        "cloth_id": 1002,
        "category": "Bottom",
        "dominant_color": "#2E4053",
        "dominant_color_name": "Navy Blue",
        "image_path": "/uploads/user_42/cloth_1002.jpg",
        "style": "Casual"
      },
      "shoe": {
        "cloth_id": 1003,
        "category": "Shoe",
        "dominant_color": "#34495E",
        "dominant_color_name": "Dark Gray",
        "image_path": "/uploads/user_42/cloth_1003.jpg",
        "style": "Casual"
      }
    },
    "compatibility_score": 0.87,
    "compatibility_breakdown": {
      "harmony_score": 0.85,
      "monochromatic_score": 0.82,
      "neutral_bonus": 0.15,
      "trend_boost": 0.05
    },
    "reasoning": "Contrast optim între alb și bleumarin. Pantofii gri neutri completează armonia cromatică. Albul și bleumarin sunt în tendințe pentru primăvară 2026.",
    "seasonality_match": 0.92,
    "created_at": "2026-05-21T10:30:00Z"
  },
  "response_time_ms": 445,
  "suggestions": [
    {
      "outfit_id": 5001,
      "score": 0.87,
      "rank": 1
    }
  ]
}
\end{lstlisting}

\textbf{Response Erori}:

\begin{lstlisting}[language=JSON]
// 400 - Utilizator nu are suficiente articole
{
  "success": false,
  "error": "Utilizatorul nu are suficiente articole pentru a genera o ținută completă",
  "error_code": "INSUFFICIENT_CLOTHES",
  "status": 400,
  "details": "Necesare: minimum 1 Top, 1 Bottom, 1 Shoe. Actual: Top=0, Bottom=2, Shoe=1"
}
\end{lstlisting}

---

\subsection{POST /web\_outfit}

\textbf{Descriere}: Generare outfit alternative din surse web (e-commerce integration).

\textbf{Request}:

\begin{lstlisting}[language=HTTP]
POST /web_outfit HTTP/1.1
Authorization: Bearer <token>
Content-Type: application/json
\end{lstlisting}

\begin{lstlisting}[language=JSON]
{
  "reference_outfit_id": 5001,
  "replace_component": "top",
  "budget_max_usd": 100,
  "retailers": ["Zalando", "H&M", "Uniqlo"]
}
\end{lstlisting}

\textbf{Response Succes (200)}:

\begin{lstlisting}[language=JSON]
{
  "success": true,
  "web_alternatives": [
    {
      "product_id": "zalando_abc123",
      "name": "Uniqlo White T-Shirt",
      "price_usd": 29.99,
      "url": "https://zalando.com/...",
      "image": "https://cdn.zalando.com/...",
      "retailer": "Zalando",
      "compatibility_score": 0.91,
      "in_stock": true
    },
    {
      "product_id": "hm_def456",
      "name": "H&M Cotton Tee White",
      "price_usd": 19.99,
      "url": "https://hm.com/...",
      "image": "https://cdn.hm.com/...",
      "retailer": "H&M",
      "compatibility_score": 0.88,
      "in_stock": true
    }
  ],
  "response_time_ms": 780
}
\end{lstlisting}

\textbf{Response Erori}:

\begin{lstlisting}[language=JSON]
// 404 - Outfit nu găsit
{
  "success": false,
  "error": "Outfit ID nu găsit",
  "error_code": "OUTFIT_NOT_FOUND",
  "status": 404
}
\end{lstlisting}

---

\section{Gestionare Istoric și Favorite}

\subsection{GET /outfits/history}

\textbf{Descriere}: Retrieval istoric ținute generate.

\textbf{Request}:

\begin{lstlisting}[language=HTTP]
GET /outfits/history?limit=20&offset=0&sort=created_at_desc&liked_only=false HTTP/1.1
Authorization: Bearer <token>
\end{lstlisting}

Query Parameters:
\begin{itemize}
    \item \texttt{limit} (Int, default 20)
    \item \texttt{offset} (Int, default 0)
    \item \texttt{sort} ("created_at\_desc" | "created_at_asc" | "score_desc")
    \item \texttt{liked\_only} (Boolean, default false)
\end{itemize}

\textbf{Response Succes (200)}:

\begin{lstlisting}[language=JSON]
{
  "success": true,
  "outfits": [
    {
      "outfit_id": 5001,
      "user_id": 42,
      "top_image": "/uploads/user_42/cloth_1001.jpg",
      "bottom_image": "/uploads/user_42/cloth_1002.jpg",
      "shoe_image": "/uploads/user_42/cloth_1003.jpg",
      "compatibility_score": 0.87,
      "liked": true,
      "created_at": "2026-05-21T10:30:00Z"
    }
  ],
  "total_count": 42,
  "page": 1,
  "per_page": 20
}
\end{lstlisting}

---

\subsection{POST /outfits/{id}/like}

\textbf{Descriere}: Salvare outfit ca favorit.

\textbf{Request}:

\begin{lstlisting}[language=HTTP]
POST /outfits/5001/like HTTP/1.1
Authorization: Bearer <token>
\end{lstlisting}

\textbf{Response Succes (200)}:

\begin{lstlisting}[language=JSON]
{
  "success": true,
  "outfit_id": 5001,
  "liked": true,
  "updated_at": "2026-05-21T10:35:00Z"
}
\end{lstlisting}

---

\subsection{DELETE /outfits/{id}}

\textbf{Descriere}: Ștergere outfit din istoric.

\textbf{Request}:

\begin{lstlisting}[language=HTTP]
DELETE /outfits/5001 HTTP/1.1
Authorization: Bearer <token>
\end{lstlisting}

\textbf{Response Succes (200)}:

\begin{lstlisting}[language=JSON]
{
  "success": true,
  "message": "Outfit șters cu succes",
  "outfit_id": 5001
}
\end{lstlisting}

---

\section{Utility Endpoints}

\subsection{GET /health}

\textbf{Descriere}: Health check pentru diagnostic backend.

\textbf{Request}:

\begin{lstlisting}[language=HTTP]
GET /health HTTP/1.1
\end{lstlisting}

\textbf{Response Succes (200)}:

\begin{lstlisting}[language=JSON]
{
  "status": "healthy",
  "timestamp": "2026-05-21T10:30:00Z",
  "version": "1.0.0",
  "database": "connected",
  "services": {
    "image_processing": "ok",
    "auth": "ok",
    "trends": "ok"
  }
}
\end{lstlisting}

---

\subsection{GET /trends/{season}}

\textbf{Descriere}: Retrieval trend colors pentru o sezon.

\textbf{Request}:

\begin{lstlisting}[language=HTTP]
GET /trends/spring HTTP/1.1
\end{lstlisting}

\textbf{Response Succes (200)}:

\begin{lstlisting}[language=JSON]
{
  "success": true,
  "season": "spring",
  "year": 2026,
  "trend_colors": [
    {
      "hex_code": "#FF6B6B",
      "name": "Coral",
      "popularity_score": 0.95,
      "rank": 1
    },
    {
      "hex_code": "#4ECDC4",
      "name": "Turquoise",
      "popularity_score": 0.92,
      "rank": 2
    }
  ]
}
\end{lstlisting}

---

\section{Rate Limiting și Throttling}

API implementează rate limiting pentru protecție:

\begin{itemize}
    \item \textbf{Limită}: 100 requests per minut per API key / IP
    \item \textbf{Rate-Limit Headers}: Include în răspuns
    \begin{itemize}
        \item \texttt{X-RateLimit-Limit}: 100
        \item \texttt{X-RateLimit-Remaining}: 87
        \item \texttt{X-RateLimit-Reset}: 1234567890 (Unix timestamp)
    \end{itemize}
    \item \textbf{Depășire limită}: HTTP 429 Too Many Requests
\end{itemize}

---

\section{Coduri HTTP și Erori}

\textbf{Coduri Succes}:
\begin{table}[h]
    \centering
    \begin{tabular}{|l|l|}
    \hline
    \textbf{Cod} & \textbf{Semnificație} \\ \hline
    200 & OK - Succes general \\ \hline
    201 & Created - Resursă creată (POST/PUT) \\ \hline
    204 & No Content - Succes fără body (DELETE) \\ \hline
    \end{tabular}
\end{table}

\textbf{Coduri Client Error}:
\begin{table}[h]
    \centering
    \begin{tabular}{|l|l|}
    \hline
    \textbf{Cod} & \textbf{Semnificație} \\ \hline
    400 & Bad Request - Invalid input \\ \hline
    401 & Unauthorized - Missing/invalid token \\ \hline
    403 & Forbidden - Permission denied \\ \hline
    404 & Not Found - Resource not found \\ \hline
    409 & Conflict - Resource already exists \\ \hline
    413 & Payload Too Large - File exceeds limit \\ \hline
    429 & Too Many Requests - Rate limit exceeded \\ \hline
    \end{tabular}
\end{table}

\textbf{Coduri Server Error}:
\begin{table}[h]
    \centering
    \begin{tabular}{|l|l|}
    \hline
    \textbf{Cod} & \textbf{Semnificație} \\ \hline
    500 & Internal Server Error - General server error \\ \hline
    503 & Service Unavailable - Backend down/overloaded \\ \hline
    \end{tabular}
\end{table}

---

\section{Exemple de Integrazione Client}

\subsection{JavaScript/React Native - Axios}

\begin{lstlisting}[language=JavaScript]
import axios from 'axios';

const api = axios.create({
  baseURL: 'https://api.stylegenai.com',
  timeout: 10000,
});

// Interceptor pentru JWT
api.interceptors.request.use(config => {
  const token = localStorage.getItem('accessToken');
  if (token) {
    config.headers.Authorization = `Bearer ${token}`;
  }
  return config;
});

// Exemplu: Login
async function login(email, password) {
  try {
    const response = await api.post('/auth/login', { email, password });
    localStorage.setItem('accessToken', response.data.access_token);
    return response.data;
  } catch (error) {
    console.error('Login failed:', error.response.data);
  }
}

// Exemplu: Upload haina
async function uploadCloth(imageFile, category) {
  const formData = new FormData();
  formData.append('image', imageFile);
  formData.append('category', category);

  try {
    const response = await api.post('/clothes/upload', formData, {
      headers: { 'Content-Type': 'multipart/form-data' },
    });
    return response.data;
  } catch (error) {
    console.error('Upload failed:', error.response.data);
  }
}

// Exemplu: Get suggestion
async function getSuggestion(season, occasion) {
  try {
    const response = await api.post('/get_suggestion', {
      season,
      occasion,
      limit_combinations: 1,
    });
    return response.data.outfit;
  } catch (error) {
    console.error('Suggestion generation failed:', error.response.data);
  }
}
\end{lstlisting}

---

\section{Testing cu cURL}

\subsection{Exemplu 1: Login}

\begin{lstlisting}[language=bash]
curl -X POST https://api.stylegenai.com/auth/login \
  -H "Content-Type: application/json" \
  -d '{
    "email": "user@example.com",
    "password": "SecurePassword123!"
  }'
\end{lstlisting}

\subsection{Exemplu 2: Upload Haina}

\begin{lstlisting}[language=bash]
curl -X POST https://api.stylegenai.com/clothes/upload \
  -H "Authorization: Bearer <token>" \
  -F "image=@/path/to/image.jpg" \
  -F "category=Top" \
  -F "style=Casual"
\end{lstlisting}

\subsection{Exemplu 3: Get Suggestion}

\begin{lstlisting}[language=bash]
curl -X POST https://api.stylegenai.com/get_suggestion \
  -H "Authorization: Bearer <token>" \
  -H "Content-Type: application/json" \
  -d '{
    "season": "spring",
    "occasion": "casual"
  }'
\end{lstlisting}

---

\section{Changelog API}

\begin{table}[h]
    \centering
    \begin{tabular}{|l|l|p{4cm}|}
    \hline
    \textbf{Versiune} & \textbf{Data} & \textbf{Modificări} \\ \hline
    1.0.0 & 2026-05-21 & MVP release cu endpoints core \\ \hline
    1.1.0 (planned) & 2026-06-15 & Adăugare /web\_outfit endpoint \\ \hline
    2.0.0 (planned) & 2026-09-01 & Reinforcement learning integration \\ \hline
    \end{tabular}
\end{table}

